\paragraph{Volume I} (the present volume) gives the text of EXI-01-a divided into numbered lemmata, a commentary on every lemma and on every declaration of the Lean module that formalizes Theorem~4, and three computed appendices: a concordance of the words of the text across the manuscripts of the discipline, a table of the numbered statements of those manuscripts with the lemmata they bear on, and the testimonia, that is, every sentence of those manuscripts that cites EXI-01. It contains the defence of Assumption~2 (Section~\ref{s:defence}) and the discussion of the absurd (Section~\ref{s:absurd}) to which Remark~5 of EXI-01-a refers.

\paragraph{Volume II, Objective Worth and the Absurd,} extends the defence of Assumption~2 given in Section~\ref{s:defence}. Part~A states the theory of value that the restricted notion of objective worth requires and shows that it is supplied by the account of the good in ETH-01 and by the naturalism of MTE-01, and that the arguments from queerness, from evolutionary debunking and from disagreement do not reach it (MTE-05, MTE-09, MTE-12). Part~B examines subjectivist theories of meaning, from Taylor's Sisyphus to Frankfurt's account of caring, and the hybrid views that answer them, and argues that the dilemma usually urged against hybrids is resolved by the reading of Definition~1 given in Section~\ref{s:def}. Part~C is a history of the absurd from Ecclesiastes to Nagel, with chapters on Tolstoy's \emph{Confession}, on Camus, on Nagel's two treatments of 1971 and 1986 and on recent work on cosmic insignificance; it develops the analysis of the vindication premise given in Section~\ref{s:absurd} and relates it to EXI-04 and EXI-06. Part~D takes up theism: it relates the additional meaning that EXI-11 assigns to a theistic framework to the argument of REL-01, and shows that the conclusion of EXI-01 is independent of both.

\paragraph{Volume III, Availability, Degree and Death,} extends the commentary on Assumption~3, on the word \enquote{extent} in Definition~1 and on Remark~5. Part~A examines the scope of the availability claim: the reading of \enquote{everyone} as every agent capable of intentional action, the lives that fall outside that scope, and the hard cases of deprivation, illness and depression that EXI-01-c answers. Part~B gives a measure of the degree of meaning: it states Definition~1 with an explicit measure on the parts of a life, proves that the propositions of Volume~I hold under it, and examines the alternatives, including the shape of a life (EXI-19) and the role of achievement (EXI-17). Part~C relates meaning to death: the badness of death (EXI-02), the undesirability of immortality (EXI-03), the fittingness of fear and regret (EXI-05) and the limit that gives a life its shape (EXI-19). Part~D contains a formal appendix, the bibliography of the three volumes and the indices.

\section*{Preface}

EXI-01-a is the shortest principal manuscript of its discipline. It occupies one page and contains, under its title, eleven sentences. It is also the manuscript to which the other manuscripts of the discipline most often refer when they need an account of what makes a life meaningful: EXI-03-a uses it to restrict the categorical desires that give life its point, EXI-04-a accepts it and argues that it leaves the absurd in place, EXI-06-b uses it to show that cosmic indifference takes nothing from meaning, and EXI-11-a and EXI-19-b take its account as their starting point. A text so short and so much used is in the condition of a classical text: each phrase carries weight, and each phrase can be read in more than one way. That is the reason for a commentary.

The form of the Commentary is the old one. The text is divided into lemmata, each quoted exactly and numbered; each lemma is followed by comment on its claim, its sources, its parallels, the readings it admits, the objections it faces and its relation to other texts. Where the text is a definition or an assumption, the comment is phrase by phrase, because in a definition every phrase constrains what is defined. Where the text is a proof, the comment is step by step. Where the text refers outward, to Wolf, Metz or Nagel, or to the Commentary itself, the comment follows the reference. We have tried throughout to explain the text and not to replace it: where we propose a reading that the text does not compel, we say so and give the alternatives, and where we find the text deficient, we say what would repair it and whether the repair is consistent with the conclusion the text defends.

That conclusion is the conclusion of its family: life is meaningful to the extent that one engages with projects of objective worth, and meaning in this sense is available to everyone. The Commentary defends it. A commentary that did not would be a different kind of work, and the reader who wants the case against the conclusion will find its strongest forms stated in the objections that follow each lemma.

\paragraph{Conventions.} Lemmata of EXI-01-a are cited by number in square brackets: \Lm{2} is the sentence of Definition~1, \Lm{2.5} the phrase \enquote{active engagement} within it. The list of lemmata is given in Section~\ref{s:text}; a program (Appendix~\ref{app:conc}) checks that every quotation in the Commentary agrees with the text. Numbered statements of EXI-01-a are cited as they are numbered there (Definition~1, Assumptions~2 and~3, Theorem~4, Remark~5); numbered statements of the Commentary are cited by section (Proposition~\ref{p:redundant}). Other manuscripts are cited by identifier, and their statements by the numbers they bear in those manuscripts. Declarations of the Lean module are cited by name in typewriter type. The text of EXI-01-a uses the spellings \enquote{defence} and \enquote{sceptical}; we keep them in quotation and use American spelling elsewhere.
